\begin{center}
\textbf{I. Erzeugung von Invarianten.}
\end{center}

Eine im allgemeinen nicht abbrechende Reihe von projektiven Invarianten
von \(f(x,dx)\) wird durch die Polaren geliefert. Wegen der Linearität
der Transformation der Differentiale folgt aus
\(f(x,dx)=g(y,dy)\) auch:
\(f(x,dx+\lambda\delta x)=g(y,dy+\lambda\delta y)\); und daraus ergeben
sich durch Entwicklung nach \(\lambda\) die für die Invarianz
charakteristischen Relationen:
\srcnumdisplay{(2)}{%
\begin{gathered}
f(dx)=g(dy),\\
f_\delta=\sum\frac{\partial f(dx)}{\partial dx}\,\delta x
       =\sum\frac{\partial g(dy)}{\partial dy}\,\delta y,\\
\vdots\\
f_{\delta^\sigma}=
\sum\frac{\partial^\sigma f(dx)}{\partial dx^\sigma}\,\delta x^\sigma
=
\sum\frac{\partial^\sigma g(dy)}{\partial dy^\sigma}\,\delta y^\sigma,\\
\vdots
\end{gathered}
}

Jede Polare gibt durch Anwendung eines Variationsprozesses \(\bdelta\)
Veranlassung zu einer nicht abbrechenden Reihe allgemeiner Invarianten:
\srcnumdisplay{(3)}{%
\begin{gathered}
\bdelta^\rho f(x,dx)=\bdelta^\rho g(y,dy),\\
\vdots\\
\bdelta^\rho f_{\delta^\sigma}
=\bdelta^\rho\sum\frac{\partial^\sigma f(dx)}{\partial dx^\sigma}\,
\delta x^\sigma
=\bdelta^\rho\sum\frac{\partial^\sigma g(dy)}{\partial dy^\sigma}\,
\delta y^\sigma,\\
\vdots\\[0.2em]
(\rho=0,1,2,\ldots)
\end{gathered}
}
Die Relationen (2) und (3) lassen vermöge (1) die Transformationsgesetze
für die Ableitungen von \(f(x,dx)\) ablesen, was aber im folgenden
nicht gebraucht wird. Aus den Invarianten (3) läßt sich nun durch
lineare Kombination die ,,Normalform der \(\rho\)ten Variation``
gewinnen, die sich für \(\rho=1\) bei Lagrange, für \(\rho=2\) bei
Riemann findet; und die dadurch charakterisiert ist, daß die gemischten
Differentiale \((\rho+1)\)ter Ordnung: \(d^\rho\delta\),
\(d^{\rho-1}\delta^2,\ldots\) eliminiert sind. Man erhält für sie:
\srcnumdisplay{(4)}{%
\begin{gathered}
\Omega_1=\delta f(dx)-d f_\delta(dx),\\
\Omega_2=\delta^2 f-\delta d f_\delta+\frac12 d^2 f_{\delta^2},\\
\vdots\\
\Omega_\rho=\delta^\rho f-\delta^{\rho-1}d f_\delta
+\frac12\delta^{\rho-2}d^2 f_{\delta^2}
+\cdots+\frac{(-1)^\rho}{\rho!}\,d^\rho f_{\delta^\rho},\\
\vdots
\end{gathered}
}

Aus den Invarianten (4) läßt sich nun durch Verallgemeinerung eines
Riemannschen Ansatzes zu Invarianten gelangen, die nur erste
Differentiale enthalten; solche Invarianten sollen als
,,Grundfunktionen`` bezeichnet werden. Ersetzt man nämlich in
\(\Omega_1\) die Differentiale durch die Differentialquotienten, so gibt
\(\Omega_1=0\) (identisch in \(\bdelta x\)) genau die Lagrangeschen
Gleichungen des zu
\[
f\!\left(\frac{dx}{dt}\right)
\]
gehörigen Variationsproblems. Ich stelle nun die invariante Bedingung
auf, daß die Richtung \((d+\lambda\delta)\) diesen Lagrangeschen
Gleichungen genügt:
\srcnumdisplay{(5)}{%
\begin{gathered}
\Omega_1(d+\lambda\delta)=
\bdelta f(dx+\lambda\delta x)
-(d+\lambda\delta)f_{\bdelta}(dx+\lambda\delta x)=0,\\
(\text{identisch in }\bdelta x),
\end{gathered}
}
woraus durch die Substitution \(\bdelta=d+\lambda\delta\) für
homogenes \(f(dx)\) noch folgt:
\srcnumdisplay{(6)}{%
(d+\lambda\delta)f(dx+\lambda\delta x)=0,
}
mit Ausnahme des Falles der Homogenität erster Ordnung. Im letzteren
Falle soll (6) als neue Bedingung zugefügt werden, da dann das
Gleichungssystem (5) nur \((n-1)\) unabhängige Bedingungen darstellt.
Setzt man in (5), bezw. (5) und (6) die Koeffizienten der nullten,
ersten und zweiten Potenz von \(\lambda\) gleich Null, so erhält man
also drei invariante Gleichungssysteme
\(\varphi=0\), \(\psi=0\), \(\chi=0\) (identisch in \(\bdelta x\)),
vermöge deren sich \(d^2x,d\delta x,\delta^2x\) durch erste Differentiale
ausdrücken lassen. Die weiteren invarianten Gleichungssysteme:
\srcnumdisplay{(7)}{%
\begin{gathered}
d^\sigma\varphi=0,\qquad d^\sigma\psi=0,\qquad
d^\sigma\chi=0,\qquad d^{\sigma-1}\delta\chi=0,\\
d^{\sigma-2}\delta^2\chi=0,\ldots,\delta^\sigma\chi=0
\qquad(\sigma=0,1,2,\ldots)
\end{gathered}
}
genügen dann gerade, um alle höheren Differentiale durch erste
auszudrücken.\srcfn{1)}{Nur die Linearform \(\sum A_i\,dx_i\) bedarf
einer besonderen Behandlung, da hier das Gleichungssystem (5) keine
zweiten Differentiale enthält.} Die Funktionen (4), verbunden mit dem
invarianten Gleichungssystem (7), führen also zu Grundfunktionen
\([\Omega_\rho]\) mit nur ersten Differentialen, wobei noch
\([\Omega_1]\) identisch verschwindet.

Ebenso läßt sich aus dem Differential \(\delta h(dx,\delta x)\) einer
Grundfunktion \(h\) vermöge (7) wieder eine Grundfunktion erhalten, die
als ,,kovariante Ableitung`` \([h^{(1)}]\) von \(h\) bezeichnet werden
soll; allgemein bezeichne \([h^{(\nu)}]\) die kovariante Ableitung von
\([h^{(\nu-1)}]\). Diese beiden Prozesse führen so zu einer doppelt
unendlichen Reihe von Grundfunktionen:
\srcnumdisplay{(8)}{%
\begin{array}{cccccc}
[\Omega_2], & [\Omega_2^{(1)}], & [\Omega_2^{(2)}] & \ldots
& [\Omega_2^{(\nu)}] & \ldots\\
\vdots &&&& \vdots&\\
{}[\Omega_\rho], & [\Omega_\rho^{(1)}], & \ldots & \ldots
& [\Omega_\rho^{(\nu)}] & \ldots\\
\vdots & \vdots &&& \vdots&
\end{array}
}

Es zeigt sich ferner, daß die linken Seiten der Gleichungen, die die
höheren Differentiale durch erste ausdrücken, den ersten Differentialen
kogredient sind (d.~h. denselben linearen Transformationen unterliegen).
Führt man also statt der höheren Differentiale diese linken Seiten
\(p,q,r,\ldots\) als Argumente der Invariante ein, so hängt diese außer
von den Ableitungen nur von zu einander kogredienten Größen ab. Die
oben angegebene Bildung der Funktionen (8) kommt einfach auf die
Einführung dieser neuen Größen hinaus; jede Invariante geht dadurch in
eine Summe von Invarianten über, und man greift daraus diejenige heraus,
die gerade \(dx,\delta x\), nicht aber \(p,q,r,\ldots\) enthält.

Im folgenden wird gezeigt, daß unter der Homogenitäts-Voraussetzung
\[
\text{--- }\;f(\varkappa\cdot dx)=\Phi(\varkappa)\cdot f(dx)\;\text{ ---}
\]
die Grundfunktionen (8) das gesuchte vollständige System erschöpfen.

